\documentclass[11pt,a4paper]{article}
\usepackage[margin=27mm]{geometry}
\usepackage{fontspec,unicode-math,amsmath,amsthm,booktabs,longtable,hyperref}
\setmainfont{STIXGeneral}\setmathfont{STIX Two Math}\hypersetup{hidelinks}
\newtheorem{theorem}{Theorem}\newtheorem{lemma}[theorem]{Lemma}
\newtheorem{proposition}[theorem]{Proposition}
\newcommand{\PSL}{\operatorname{PSL}}\newcommand{\Aut}{\operatorname{Aut}}
\title{The universal dodecahedral--hemi-icosahedral polytope}
\author{A complete corrected recovery of Hartley--Leemans}
\date{}
\begin{document}\maketitle
Hartley and Leemans determine a finite universal locally projective polytope and its quotients \cite{source}. Their later erratum removes eight quotient classes \cite{erratum}. We recover the presented group and both distinguished quotients, then apply the separately reviewed exhaustive census. The finite image alone is not an order proof: the missing upper bound is supplied by coset enumeration and a relation calculation in a cyclic subgroup. Every program used for that step is printed in the accompanying appendix.

The source here is the complete fifteen-page 2003 arXiv edition, not an assumed copy of the twelve-page journal article. The complete census paper, tables, programs and accepted computation are included in the declared dependency packet, without repeating the subgroup enumeration. They are part of the proof boundary and must be charged in any later measurement.

\section{Finite presentations and matching bounds}
\begin{lemma}\label{bound}
Let $W$ be a presented group, let $L\leq W$ have finite index $d$, and suppose that a presented group $T$ maps onto $L$. Suppose $T$ has a cyclic subgroup $C=\langle x\rangle$ of index $e$. If a complete integer relation matrix for $C/[C,C]$, on the specified generator $x$, has one column with nonzero gcd $m$, then
\[
 |W|\leq dem.
\]
If $W$ maps onto a finite marked group of order $dem$, that epimorphism is an isomorphism.
\end{lemma}
\begin{proof}
A group generated by one element is abelian. Thus $C=C/[C,C]\cong\mathbb Z/m\mathbb Z$: its defining relation subgroup in $\mathbb Z$ is generated by the matrix entries. Hence $|T|=em$, while surjectivity gives $|L|\leq em$. Coset partitioning now gives $|W|=d|L|\leq dem$. A surjection onto a group of this upper-bound order has trivial kernel.
\end{proof}
The hypothesis concerns a complete subgroup relation matrix, not a matrix from an arbitrary finite quotient. In our calculation it is computed by GAP's modified Todd--Coxeter algorithm for the subgroup generated by the actual element $x$ of the presented group. The simpler finite permutation image is used only on the opposite side of the inequality. GAP's exact presentation algorithms and ACE's completed coset enumeration are retained computational premises; neither a list of permutations satisfying the relators nor a closed table with no verified subgroup enumeration would replace them.

\begin{lemma}\label{prime}
Let $K$ be finite and let $B\subseteq K$. Suppose that for every prime $p\mid |K|$ and every conjugacy class of subgroups of order $p$, some member meets $B$. Then every nontrivial subgroup $N\leq K$ has a conjugate meeting $B$.
\end{lemma}
\begin{proof}
Cauchy's theorem gives a subgroup $C\leq N$ of prime order. By hypothesis $C^g\cap B\ne\varnothing$ for some $g$, and this intersection is contained in $N^g\cap B$.
\end{proof}
There is a complete finite test without enumerating every subgroup: form a cyclic subgroup from each element-class representative of order $p$, identify these subgroups up to conjugacy, and test their conjugates against the elements of $B$ of order $p$. Every subgroup of order $p$ occurs because any of its nonidentity elements generates it. This proof does not require all elements of order $p$ to be conjugate.

\section{The literal universal group}
Let $W$ have generators $s_0,s_1,s_2,s_3$ and relations
\begin{equation}\label{W}
 s_i^2=1,\quad(s_0s_1)^5=(s_1s_2)^3=(s_2s_3)^5=1,
 \quad[s_i,s_j]=1\ (|i-j|>1),\quad(s_1s_2s_3)^5=1.
\end{equation}
Put $\omega=(s_0s_1s_2)^5$ and $\nu=\omega s_3$.
\begin{theorem}\label{structure}
The literal presented group is marked-isomorphic to
\[
 W\cong J_1\times\PSL_2(19),\qquad |W|=600415200.
\]
The normal closures $N'=\langle\!\langle\nu^3\rangle\!\rangle$ and $N''=\langle\!\langle\omega\rangle\!\rangle$ are respectively the $\PSL_2(19)$ and $J_1$ factors. In particular $W/N'\cong J_1$, $W/N''\cong\PSL_2(19)$, and $\nu$ has order six.
\end{theorem}
\begin{proof}
Use exactly the source subgroup
\[
 L=\langle u_1,u_2,u_3\rangle,\quad
 u_1=s_2s_1,\quad u_2=(s_0s_1)^2s_2s_3s_0,\quad
 u_3=s_0s_1s_0s_3s_2s_3s_1.
\]
The complete ACE enumeration on presentation \eqref{W} and these three subgroup words has index $20$. Reduced Reidemeister--Schreier rewriting followed by tracked Tietze transformations produces two generators $X,Y$ and the exact relators
\begin{equation}\label{T}
 Y^3=1,\qquad XYXY^{-1}XYX^{-1}YX=1.
\end{equation}
The tracked images of the three primary subgroup generators are
\[
 u_1=X^{-1}Y^{-1}X,\quad u_2=Y^{-1}X^{-1}Y,
 \quad u_3=Y^{-1}X^{-1}Y^2X.
\]
Consequently the source words $x=u_2^{-1}u_3u_1$ and $y=u_1u_3^{-1}$ become $X$ and $X^{-1}Y^{-3}XY=Y$. This checks the marked source presentation, not merely an abstract presentation of some subgroup. Both the complete initial presentation and all tracked generator images are saved.

In the group $T$ defined by \eqref{T}, modified Todd--Coxeter enumeration for the specified subgroup $\langle X\rangle$ gives index $175560$. Its one-column abelianized subgroup relation matrix has $290$ nonzero rows, all divisible by $171$, and contains the row $-171$. Thus its gcd is exactly $171$. Lemma~\ref{bound} gives
\[
 |W|\leq20\cdot175560\cdot171=600415200.
\]
The complete code performs these calculations in the presented groups before introducing any finite image.

For the lower bound we reuse the complete Janko marking search, already printed and checked in the census packet. It constructs four elements $j_i$ in the named TomLib model of $J_1$. The other four markers $p_i$ are the literal permutations on twenty points printed on source page six. The exact finite checks show
\[
 \langle(j_i,p_i):0\leq i\leq3\rangle=J_1\times P,
 \quad |J_1|=175560,\quad P\cong\PSL_2(19),\quad |P|=3420.
\]
Every relation in \eqref{W} is checked in this product. The named $J_1$ model is retained library data; the twenty-point group is additionally compared by an actual isomorphism with GAP's $\PSL_2(19)$ construction. The resulting epimorphism therefore has the upper-bound order and is an isomorphism.

In the Janko coordinates the images of $\omega,\nu$ have orders $2,3$; in the projective coordinates they have orders $1,2$. Thus $\omega$ belongs to the Janko factor and $\nu^3$ belongs to the projective factor, both nontrivially. Simplicity of the factors, also checked in the finite models, makes their normal closures the entire indicated factors. This proves the marked quotient presentations, since adding the respective word is exactly quotienting by its normal closure. The two coordinate orders of $\nu$ give order six.
\end{proof}

The source's twenty-coset action is retained literally in the data and has kernel $N''$. Our direct-product construction supplies the faithful action on $266+20=286$ points described in the source; it is not claimed to reproduce the unavailable historical permutation file. The preceding proof also recovers the source's two-generator presentation of $L$, its order $30020760$, and the composition factors. It avoids any reliance on that file, \texttt{g.txt}.

\section{Geometry and the two distinguished quotients}
We use the classical marked Coxeter group $[5,3]$ of order $120$ for the dodecahedron, its reversed marking for the icosahedron, and the antipodal quotient of order $60$ for the corresponding hemi-polyhedron. The central antipode is the fifth power of the product of the three distinguished reflections. We also use the standard string-C-group/polytope correspondence and the general semisparse quotient correspondence. These are retained geometric inputs, not the target universal-group theorem.

\begin{theorem}\label{geometry}
The group in Theorem~\ref{structure}, with its displayed marking, is a string C-group. Its regular polytope $\mathcal P$ has dodecahedral facets and hemi-icosahedral vertex figures. It is the universal polytope with these local types. The quotient with group $J_1$ has $1463$ dodecahedral facets and $2926$ vertices; the quotient with group $\PSL_2(19)$ is the self-dual $57$-cell with hemi-dodecahedral facets and hemi-icosahedral vertex figures. Both distinguished quotients have no proper polytope quotient under the semisparse convention.
\end{theorem}
\begin{proof}
For each of the three actual markings, the required program constructs the subgroup on every subset of the four markers and checks all $16^2=256$ intersection equalities. The markers are nonidentity involutions, their adjacent products have exact orders $5,3,5$, and the nonadjacent pairs commute. These checks establish the entire string C-property, not only its two endpoint intersections.

The endpoint orders are respectively $(120,60)$ for $W$ and $J_1$, and $(60,60)$ for the projective quotient, where the first entry is the facet group. Their Coxeter presentations give surjections from the classical local groups just described, and equality of their orders makes them marked isomorphisms. Thus their local polytopes are the claimed ones; an abstract group isomorphism to $C_2\times A_5$ alone would not suffice. The string-C-group correspondence constructs the regular polytopes. Group order divided by the appropriate endpoint order gives $1463$, $2926$ and $57$. Reversing the four actual projective markers extends to a bijective group homomorphism, explicitly checked by the program, and hence proves the asserted self-duality.

For universality, take any polytope with the indicated local types and let its involutory flag adjacencies be $r_i$. Nonadjacent adjacencies commute; the polygon sections give the three adjacent relations. At each vertex the hemi-icosahedral flag action gives $(r_1r_2r_3)^5=1$. Thus \eqref{W} acts transitively on the flags, and the stabilizer of a chosen flag identifies the polytope with the corresponding quotient of $\mathcal P$ by the general quotient correspondence. This applies without assuming the target regular. It also applies when some facets have been replaced by their antipodal quotients, since their adjacencies still satisfy the same facet relations. Every section-regular locally projective polytope of type $\{5,3,5\}$ is of one of these local types or has dual of one of them. This proves the precise classification coverage claimed in the source; the mixed-facet quotients below are additional, non-section-regular members of the full quotient census.

For the final assertion retain the classical fact that the hemi-icosahedron has no proper quotient and Hartley's general endpoint semisparseness theorem. Applied to these local groups, it says that a subgroup $N$ of the Janko marked group is semisparse exactly when
\[
 N^g\cap H_0H_3\subseteq\{1,\omega\}\quad\hbox{for all }g;
\]
in the projective quotient the right side is $\{1\}$. The facet alternatives are respectively the dodecahedron and its antipodal quotient, or only the hemi-dodecahedron. For each prime and every conjugacy class of prime-order subgroups, the finite program finds a member meeting the corresponding forbidden endpoint product. Lemma~\ref{prime} therefore excludes every nontrivial semisparse subgroup.

This calculation also checks the source's explicit Janko witnesses: in its marker indices the words
\[
 s_1s_2,\quad s_0s_1,\quad s_3s_1s_2s_1s_0,\quad
 s_3s_2s_1s_0s_1s_0,\quad
 s_1s_3s_2s_1s_3s_0s_1s_0s_2s_1s_0
\]
belong to $H_0H_3$ and have orders $3,5,19,11,7$. Its unique involution class and its prime-order odd Sylow groups give the original short Cauchy argument as well. For the projective quotient the orders tested are $2,3,5,19$; the test is on prime-order subgroups, so does not incorrectly require a Sylow $3$-subgroup to have order three.
\end{proof}

\section{The corrected full quotient classification}
\begin{theorem}\label{census}
The full quotient classification for $\mathcal P$ has $137$ isomorphism classes, with $26$ maximal semisparse subgroup classes and $349$ maximal-containment edges. It contains $65$ quotients with the original local types, one with only hemi-dodecahedral facets, and $71$ with mixed facets. Exactly $66$ are section regular. All facet counts, all actual subgroup representatives and all quotient automorphism groups are specified in the included complete census tables.
\end{theorem}
\begin{proof}
Theorem~\ref{structure} discharges the marked universal-group input that was retained in the earlier erratum recovery. The local identification in Theorem~\ref{geometry} supplies its geometric hypotheses. Hartley's endpoint criterion becomes
\[
 N^w\cap H_0H_3\subseteq\{1,\omega\}\quad(w\in W).
\]
The included census paper proves the conjugacy-class avoidance test, the facet orbit formula and the maximality criterion. Its exhaustive required computation enumerates all $1262$ subgroup classes of this exact marked finite group, filters them to $137$, and supplies the $349$ edges and $26$ maximal classes. It is a complete computation, not a sample or a retained erroneous source total. The accepted literal-row reconciliation checks all $145$ printed rows, keeps the $137$ survivors, and exhibits actual forbidden conjugacy witnesses for the eight deletions
\[
 13,70,78,90,91,95,111,127.
\]
The thirty literal source maximal-subgroup words provide $26$ surviving anchors; the other four deleted labels are fixed by their printed parent/child relations and exact local invariants. The orbit-count proof and normalizer quotient formula give all facet and automorphism data. The full checked packet, including every generated subgroup and table, is bound unchanged; no second census is asserted or needed here.
\end{proof}

\paragraph{What the source tables do and do not specify.}
The classification gives exact permutation subgroups, stronger data than semidirect-product abbreviations. The historical subgroup-description column is checked for order, while the full printed automorphism column is checked up to the elementary direct-product identities described in the census paper. We do not infer an abstract subgroup isomorphism merely from its abbreviated type and order. In particular source rows $38,39,40$ all use $19:9$-style notation, while the actual centers have orders $1,1,3$; the distinct action at row $40$ remains explicit. Exact representatives preserve all actions.

The accessible arXiv edition has these additional discrepancies, kept distinct from the erratum and from the unseen journal edition. Table~1's order-$70$ row labelled $58$ corresponds to row $63$ in the complete tables; their row $58$ instead has order $76$. Page~14 row $21$ prints $2^2\times(19:3)$, whereas its order is $456$ and its actual group is $2^3\times(19:3)$. Row $27$ prints $2\times19:3$, whereas its order is $342$ and its actual group is $2\times(19:9)$. Row $41$ has an elementary abelian normal subgroup of order eight and exponent two; its printed $8:(7:3)$ is not literally a cyclic-eight semidirect product under the source's stated notation. These exact finite distinctions, and the $38$--$40$ center distinction, are recorded separately. None is silently transferred to the journal text.

The source's numerical timing and former software history are preserved as historical reports, not mathematical results to reproduce. Its $145$, $30$, $69$, $70$ and $75$ census totals are replaced, respectively, by $137$, $26$, $65$, $66$ and $71$ in their stated categories. Its two named simple quotients, universal presentation, local types, covering scope and quotient invariants are all retained in the corrected principal contract.

\paragraph{Computation, reuse and limits.}
The source supplied the subgroup words, two-generator presentation, twenty-point marking and the principal group and geometric results. The Janko marking search is previously reviewed work, copied in full. Our new completion is an explicit bounded presentation upper bound, its cyclic-subgroup reduction, and the complete marked/source adapters. The already reviewed L084 conjugacy census is reused once in full; its conditional use of the universal theorem is discharged here, so it is not a circular premise for Theorem~\ref{structure}. Generic quotient theory, classical local Coxeter facts and named library data remain declared inputs. All three new or reused GAP files are printed, with exact runtime/data pins. A fresh replay checks $768$ subset intersection identities, both prime obstructions, every defining relator, normal closures and duality. No end-to-end Lean proof or new formal declaration is claimed, and no conjecture or historical-priority claim is made.

\begin{thebibliography}{9}
\bibitem{source} M. I. Hartley and D. Leemans, \emph{Quotients of a Universal Locally Projective Polytope of type $\{5,3,5\}$}, arXiv:math/0307075v1, 5 July 2003, fifteen pages. Journal reference: Math. Z. \textbf{247} (2004), 663--674; that edition has not been collated here.
\bibitem{erratum} M. I. Hartley and D. Leemans, erratum, Math. Z. \textbf{253} (2006), 433--434, doi:10.1007/s00209-005-0872-z.
\bibitem{census} \emph{Conjugacy tests for quotients of a locally projective polytope}; \emph{Corrected quotient census}; complete required programs, frozen L084 recovery and distinct root review. Included byte-exact in this packet.
\bibitem{arp} P. McMullen and E. Schulte, \emph{Abstract Regular Polytopes}, Cambridge University Press, 2002. Retained classical local groups and group/polytope correspondence.
\bibitem{hartley} M. I. Hartley, \emph{Simpler Conditions for Semisparse Subgroups}, source reference~[7]. Its general endpoint criterion is retained through the precise statement on source page three; no new whole-source recovery or fresh primary reading is claimed.
\end{thebibliography}
\end{document}
